\chapter{Discussion}
\label{chap:discussion}

\section{Validation des biomarqueurs}

Cohérence avec littérature : ESCO2 et ANP32A-IT1 confirment rôles immunitaire et cellulaire \cite{Zhang2022}.
Pertinence pour populations africaines.

\section{Comparaison avec l'état de l'art}

Performances comparables/supérieures aux modèles occidentaux malgré dataset synthétique.
Avantage du focus populations locales.

\section{Limitations de l'étude}

\subsection{Données synthétiques}
Validation sur données réelles camerounaises nécessaire.

\subsection{Taille du dataset}
1000 patients : suffisant pour proof-of-concept, mais échelle clinique requiert plus.

\subsection{Biomarqueurs limités}
Seulement 3 biomarqueurs testés. Autres candidats possibles.

\section{Implications pratiques}

Outil d'aide à la décision pour cliniciens camerounais.
Accessibilité (pas de GPU requis), interface intuitive.

\section{Perspectives}

Validation externe, intégration autres biomarqueurs, études longitudinales.
